%%%%%%%%%%%%%%%%%%%%%%%%%%%%%%%%%%%%%%%%%%%%%%%%%%%%%%%%%%%%%%%%%%%%%%%%%%%%%
%% sn-article.tex — SN Computer Science formatted version
%% Stretching the Manifold paper
%%%%%%%%%%%%%%%%%%%%%%%%%%%%%%%%%%%%%%%%%%%%%%%%%%%%%%%%%%%%%%%%%%%%%%%%%%%%%

\documentclass[pdflatex,sn-mathphys-num]{sn-jnl}

% ── Extra packages used by your content ──────────────────────────────────────
\usepackage{amsmath,amssymb,amsthm}
\usepackage{graphicx}
\usepackage{booktabs}

% ── Theorem-like environments ────────────────────────────────────────────────
% sn-jnl already provides \theoremstyle and some environments via amsthm.
% We define ours with unique names to avoid clashes.
\theoremstyle{thmstyleone}
\newtheorem{definition}{Definition}
\newtheorem{proposition}{Proposition}
\newtheorem{observation}{Observation}

\raggedbottom

\begin{document}

\title[Stretching the Manifold]{Stretching the Manifold: The Ambient-Intrinsic Dimension Gap as a Universal Predictor of Machine Learning Algorithm Behavior}

\author*[1]{\fnm{Dayanch} \sur{Amanov}}\email{amanovvdayanch@gmail.com}

\affil*[1]{\orgname{China University of Petroleum (Beijing)}, \orgaddress{\street{No. 18 Fuxue Road}, \city{Beijing}, \postcode{102249}, \country{China}}}

\abstract{High-dimensional data is routinely analyzed using algorithms designed for
low-dimensional settings, yet predicting when such algorithms will fail remains
an open problem. Prior work has studied failure as a function of ambient
dimension or intrinsic dimension separately, but neither alone provides a
reliable predictor. We propose the stretching ratio $\rho = d_a / d_i$
-- the ratio of ambient to intrinsic dimension -- as a unified predictor of
algorithm behavior. Through controlled experiments across 45 conditions, we
show that $\rho$ follows a power law relationship with k-means clustering
failure, with a universal phase transition at $\rho^* \approx 3.5$--$5$
independent of intrinsic dimension. The decay exponent $\beta$ scales
logarithmically with $d_i$ as $\beta(d_i) = 0.140 \ln(d_i) + 0.534$.
Strikingly, $\rho$ has the opposite effect on distance-weighted local
methods: kNN classification and regression improve monotonically with $\rho$,
revealing a geometric tradeoff rather than a simple curse. We further show that
$\log(\rho)$ explains 55.6\% of ARI variance, outperforming participation ratio
(32.1\%), explained variance ratio (31.3\%), and the ambient-intrinsic baseline
(14.7\%) as individual predictors. Under heavy-tailed noise, the failure
threshold shifts to $\rho^* < 2$, connecting our framework to prior work on
distributional effects in high-dimensional spaces. Our results provide
actionable thresholds for practitioners: measure $\rho$ before running any
clustering pipeline, and treat $\rho > 5$ as a danger zone.}

\keywords{Curse of dimensionality, Intrinsic dimension, Stretching ratio, k-means clustering, Manifold hypothesis, Scaling laws}

\maketitle

\section{Introduction}

The \textit{curse of dimensionality} --- the degradation of algorithm 
performance as data dimensionality increases --- is one of the most studied 
phenomena in machine learning. Yet despite decades of research, practitioners 
lack simple, computable thresholds that predict \textit{when} a given algorithm 
will fail on their specific dataset.

Prior work has approached this problem from two directions. The first studies 
failure as a function of \textit{ambient dimension} $d_a$ --- the raw number of 
features. The second, motivated by the manifold hypothesis, focuses on 
\textit{intrinsic dimension} $d_i$ --- the true degrees of freedom of the data 
distribution. Both capture important aspects of complexity, but neither alone 
reliably predicts failure across diverse settings.

We propose a different perspective: what matters is not $d_a or d_i$ 
individually, but the \textit{gap between them} --- how far the data manifold 
is ``stretched'' within the ambient space. Formally, we define the 
\textit{stretching ratio}:
\begin{equation}
    \rho = \frac{d_a}{d_i}
\end{equation}
When $\rho = 1$, the data fills its ambient space efficiently. As $\rho$ grows, 
an increasing proportion of ambient dimensions carry no signal --- only noise 
--- diluting pairwise distances and distorting geometric structure.

\paragraph{Contributions.} We make the following contributions:
\begin{itemize}
    \item We show that k-means clustering failure follows a power law in $\rho$: 
    $\text{ARI}(\rho) \approx \alpha \cdot \rho^{-\beta}$, with a universal 
    phase transition at $\rho^* \approx 3.5$--$5$.
    \item We derive an empirical law for the decay exponent: 
    $\beta(d_i) = 0.140 \ln(d_i) + 0.534$, with prediction error under 0.025.
    \item We show that kNN methods exhibit the \textit{opposite} trend --- 
    improving with $\rho$ --- revealing a fundamental algorithm-dependent 
    geometric tradeoff.
    \item We demonstrate that $\log(\rho)$ outperforms existing complexity 
    measures including participation ratio and explained variance ratio as a 
    predictor of clustering failure ($R^2=0.556$ vs 0.321).
    \item We show that heavy-tailed noise shifts $\rho^*$ downward, connecting 
    our framework to distributional failure regimes.
\end{itemize}

% ── 2. Related Work ───────────────────────────────────────────────────────────
\section{Related Work}

\paragraph{Curse of dimensionality.}
The degradation of distance-based methods in high dimensions is well 
documented~\cite{bellman1957}. Beyer et al.~\cite{beyer1999} showed that 
under certain conditions, nearest-neighbor queries become meaningless as 
dimension grows. Subsequent work characterized the rate of distance 
concentration for specific distributions.

\paragraph{Intrinsic dimensionality.}
The manifold hypothesis~\cite{bengio2013} posits that real data lies on 
low-dimensional manifolds embedded in high-dimensional space. Estimators for 
intrinsic dimension include the TwoNN method~\cite{facco2017}, maximum 
likelihood estimation~\cite{levina2004}, and correlation dimension. Our work 
uses TwoNN as the primary estimator.

\paragraph{Distance concentration and scaling laws.}
Recent work has studied how distance contrast degrades with dimensionality, 
finding consistent empirical laws across distributions with finite third 
moment. We extend this line of inquiry by studying the 
\textit{ratio} of ambient to intrinsic dimension as the governing quantity, 
and by examining failure across multiple algorithm families.

\paragraph{Adversarial vulnerability and dimension gap.}
The ambient-intrinsic gap has been studied in the specific context of 
adversarial vulnerability, where higher gap correlates 
with increased susceptibility to adversarial perturbations. We generalize this 
finding to arbitrary algorithm failure across clustering, classification, and 
regression.

% ── 3. Theoretical Framework ──────────────────────────────────────────────────
\section{Theoretical Framework}

\begin{definition}[Stretching Ratio]
Let $d_a \in \mathbb{N}$ denote the ambient dimension of a dataset 
$\mathbf{X} \in \mathbb{R}^{n \times d_a}$, and let $d_i$ denote its intrinsic 
dimension estimated via TwoNN. The \textit{stretching ratio} is defined as:
\begin{equation}
    \rho = \frac{d_a}{d_i}, \quad \rho \geq 1
\end{equation}
\end{definition}

\begin{definition}[Stretching Gap]
The \textit{stretching gap} is $\delta = d_a - d_i$, measuring the absolute 
number of noise dimensions.
\end{definition}

\paragraph{Intuition.}
When $\rho$ is large, $d_a - d_i$ dimensions contribute noise orthogonal to 
the signal manifold. Under independent noise, each noise dimension adds 
variance to pairwise distances independently, diluting the relative contrast 
between inter- and intra-cluster distances. This mechanism connects directly to 
distance concentration theory: the signal-to-noise ratio in pairwise distances 
scales as $O(d_i / d_a) = O(1/\rho)$.

\begin{proposition}[Distance Contrast Degradation]
\label{prop:contrast}
Let $\mathbf{X} \in \mathbb{R}^{n \times d_a}$ be a dataset whose signal lies on a
$d_i$-dimensional subspace, with isotropic Gaussian noise
$\epsilon \sim \mathcal{N}(0, \sigma^2 \mathbf{I}_{d_a})$ in all ambient dimensions.
Define the \textit{distance contrast ratio} as
$\Delta(\rho) = \mathbb{E}[\|\mathbf{x}_i - \mathbf{x}_j\|] / \mathrm{std}[\|\mathbf{x}_i - \mathbf{x}_j\|]$,
where the expectation is over pairs drawn from different clusters.
Then $\Delta(\rho)$ degrades as $\Omega(\rho^{-1/2})$:
\begin{equation}
    \Delta(\rho) \;=\; \Omega\!\left(\rho^{-1/2}\right).
\end{equation}
\end{proposition}

\begin{proof}[Proof sketch]
Decompose pairwise squared distance as
$\|\mathbf{x}_i - \mathbf{x}_j\|^2 = \|\mathbf{s}_i - \mathbf{s}_j\|^2 + \|\boldsymbol{\epsilon}_i - \boldsymbol{\epsilon}_j\|^2$,
where $\mathbf{s}_k$ is the signal component.
The noise term is $\chi^2$-distributed with $d_a$ degrees of freedom, so
$\mathbb{E}[\|\boldsymbol{\epsilon}_i-\boldsymbol{\epsilon}_j\|^2] = 2\sigma^2 d_a$
and $\mathrm{Var}[\|\boldsymbol{\epsilon}_i-\boldsymbol{\epsilon}_j\|^2] = 8\sigma^4 d_a$.
By concentration of measure~\cite{vershynin2018}, the noise term concentrates
around its mean at rate $O(d_a^{-1/2})$, while the signal difference
$\|\mathbf{s}_i - \mathbf{s}_j\|$ is bounded below by the cluster separation constant
in $d_i$ dimensions.
Hence the signal-to-noise ratio in Euclidean distance scales as
$O\!\left(\sqrt{d_i / d_a}\right) = O\!\left(\rho^{-1/2}\right)$,
giving the stated bound.
This is consistent with the empirically observed power law
$\mathrm{ARI}(\rho) \approx \alpha \cdot \rho^{-\beta}$ (Section~\ref{sec:results}),
with the theoretical lower bound corresponding to $\beta = 0.5$.
Observed $\beta > 0.5$ at higher $d_i$ reflects compounding effects at the
centroid estimation step of k-means; a tighter analysis is left as future work.
\end{proof}

% ── 4. Experimental Setup ─────────────────────────────────────────────────────
\section{Experimental Setup}

\subsection{Synthetic Data Generation}
We generate controlled datasets by embedding $d_i$-dimensional Gaussian data 
into $d_a$-dimensional ambient space via random projection, followed by 
isotropic Gaussian noise:
\begin{equation}
    \mathbf{X} = \mathbf{X}_{\text{low}} \mathbf{P} + \epsilon, \quad 
    \mathbf{X}_{\text{low}} \in \mathbb{R}^{n \times d_i}, \quad 
    \mathbf{P} \in \mathbb{R}^{d_i \times d_a}, \quad 
    \epsilon \sim \mathcal{N}(0, \sigma^2 \mathbf{I})
\end{equation}
where columns of $\mathbf{P}$ are normalized. We vary $d_i \in \{2,3,5,8,10\}$ 
and $\rho \in \{1,2,5,10,20,30,50,75,100\}$, yielding 45 experimental 
conditions. Each condition is repeated over 5 trials with different random seeds.

\subsection{Algorithms}
We evaluate three algorithm families:
\begin{itemize}
    \item \textbf{Clustering:} k-means with $k=3$, measured by Adjusted Rand 
    Index (ARI). Failure defined as ARI $< 0.5$.
    \item \textbf{Classification:} $k$-nearest neighbor classifier ($k=5$), 
    measured by 3-fold cross-validation accuracy.
    \item \textbf{Regression:} $k$-nearest neighbor regressor ($k=5$), measured 
    by 3-fold cross-validation $R^2$.
\end{itemize}

\subsection{Intrinsic Dimension Estimation}
We use the TwoNN estimator~\cite{facco2017}, which estimates intrinsic 
dimension from the ratio of distances to the first and second nearest neighbors. 
We validate robustness to estimator choice using MLE estimation as a secondary 
check.

\subsection{Baseline Measures}
We compare $\rho$ against: ambient dimension $d_a$, intrinsic dimension $d_i$, 
participation ratio (PR), explained variance ratio at 50\% threshold (EVR-50), 
and log condition number of the covariance matrix.

% ── 5. Results ────────────────────────────────────────────────────────────────
\section{Results}
\label{sec:results}

\subsection{Power Law Decay and Universal Failure Threshold}

Figure~\ref{fig:main} shows k-means ARI as a function of $\rho$ across all 
intrinsic dimensions. ARI degrades monotonically and follows a consistent power 
law:
\begin{equation}
    \text{ARI}(\rho, d_i) \approx \alpha \cdot \rho^{-\beta(d_i)}
\end{equation}
The phase transition --- where failure rate exceeds 50\% --- occurs consistently 
at $\rho^* \approx 3.5$--$5$ across all tested intrinsic dimensions, 
establishing a practical threshold independent of $d_i$.

\subsection{Logarithmic Scaling of the Decay Exponent}

The decay exponent $\beta$ is not constant but scales logarithmically with 
intrinsic dimension:
\begin{equation}
    \beta(d_i) = 0.140 \cdot \ln(d_i) + 0.534
\end{equation}
Prediction error is below 0.025 across all tested values (Table~\ref{tab:beta}). 
This provides a closed-form, two-parameter prediction of clustering failure 
given only $\rho$ and $d_i$.

\begin{table}[h]
\centering
\caption{Fitted decay exponents and prediction errors.}
\label{tab:beta}
\begin{tabular}{ccccc}
\toprule
$d_i$ & $\alpha$ & $\beta$ (fitted) & $\beta$ (predicted) & Error \\
\midrule
2  & 1.270 & 0.607 & 0.631 & 0.024 \\
3  & 1.311 & 0.711 & 0.687 & 0.024 \\
5  & 1.421 & 0.777 & 0.759 & 0.018 \\
8  & 1.403 & 0.813 & 0.824 & 0.011 \\
10 & 1.456 & 0.849 & 0.856 & 0.007 \\
\bottomrule
\end{tabular}
\end{table}

\subsection{Opposite Effects Across Algorithm Families}

Figure~\ref{fig:algos} reveals a striking divergence: while k-means degrades 
with $\rho$ (Spearman $r = -0.913$, $p = 2.65 \times 10^{-18}$), kNN 
classification and regression \textit{improve} (Spearman $r = +0.608$ and 
$+0.600$, both $p < 10^{-5}$). This suggests $\rho$ characterizes a 
fundamental geometric tradeoff: high $\rho$ destroys global centroid structure 
while smoothing local neighborhood geometry, benefiting distance-weighted 
local methods.

\subsection{Stretching Ratio Outperforms Existing Complexity Measures}

Table~\ref{tab:predictors} compares predictors of k-means failure. $\log(\rho)$ 
achieves $R^2=0.556$, outperforming all individual measures. Notably, 
$d_i$ alone is essentially uninformative ($R^2=0.032$), confirming that intrinsic 
dimension without ambient context carries little predictive value.

To assess whether combining predictors adds value, we also evaluate three
composite models estimated via ordinary least squares regression on the same
45-condition dataset.
(i)~$\log(\rho) + \text{PR}$ ($R^2=0.590$),
(ii)~$\log(\rho) + \text{PR} + \text{EVR-50}$ ($R^2=0.605$), and
(iii)~all seven predictors jointly ($R^2=0.625$).
Marginal gains beyond the single-predictor $\log(\rho)$ are modest,
supporting its use as a lightweight practical diagnostic.

\begin{table}[h]
\centering
\caption{Predictor comparison for k-means ARI. Higher $R^2$ is better.}
\label{tab:predictors}
\begin{tabular}{lccc}
\toprule
Predictor & Spearman $r$ & $p$-value & $R^2$ \\
\midrule
$d_a$ (ambient dim)       & $-0.728$ & $1.43 \times 10^{-8}$  & 0.147 \\
$d_i$ (intrinsic dim)     & $-0.501$ & $4.57 \times 10^{-4}$  & 0.032 \\
$\rho$ (stretching ratio) & $-0.542$ & $1.20 \times 10^{-4}$  & 0.203 \\
$\log(\rho)$              & $-0.542$ & $1.20 \times 10^{-4}$  & \textbf{0.556} \\
Participation Ratio       & $-0.856$ & $6.46 \times 10^{-14}$ & 0.321 \\
EVR-50                    & $-0.848$ & $2.08 \times 10^{-13}$ & 0.313 \\
log(Condition Number)     & $-0.639$ & $2.31 \times 10^{-6}$  & 0.351 \\
\bottomrule
\end{tabular}
\end{table}

\subsection{Heavy-Tailed Noise Shifts the Failure Threshold}

Under heavy-tailed noise ($t$-distribution, df=2), baseline ARI at $\rho=1$ 
drops by up to 36\% relative to Gaussian noise, and the failure threshold shifts 
to $\rho^* < 2$. This demonstrates that the stretching ratio interacts 
multiplicatively with distributional tail behavior, consistent with findings 
on distance concentration under non-Gaussian 
distributions~\cite{aggarwal2001}.

\subsection{Real-World Validation}

To confirm that the stretching-ratio effect generalises beyond synthetic data,
we evaluate on three real-world datasets with naturally varying $\rho$.

\paragraph{Wine (UCI).}
The Wine dataset contains 178 samples across 3 cultivar classes in 13
continuous features. We augment it by appending blocks of pure Gaussian noise
dimensions to vary $\rho$ while preserving all original cluster structure.
As shown in Table~\ref{tab:realdata}, ARI degrades sharply once $\rho$ exceeds
the synthetic threshold: from 0.897 at $\rho \approx 1.65$ to 0.208 at
$\rho \approx 3.36$, consistent with $\rho^* \approx 3.5$--$5$.

\paragraph{MNIST.}
We apply PCA to reduce the 784-dimensional MNIST embeddings to varying
$d_a \in \{10, 20, 50, 100, 200, 784\}$, estimating $d_i$ via TwoNN on a
500-sample subsample at each level. ARI remains approximately constant across
all tested $\rho$ values (Table~\ref{tab:realdata}).
This indicates that when clustering failure is dominated by non-convex cluster
geometry rather than noise-induced distance dilution, the stretching-ratio effect
is masked --- an important boundary condition of the framework.

\paragraph{Olivetti Faces.}
The Olivetti dataset (400 images, 40 subjects, 4096 features) exhibits a
similar pattern to MNIST: performance is bounded by class overlap in
face-space rather than ambient dimensionality.

\begin{table}[h]
\centering
\caption{Real-world validation results. $d_a$ and $d_i$ are estimated values; $\dagger$ noise dimensions appended.}
\label{tab:realdata}
\begin{tabular}{llcccc}
\toprule
Dataset & Condition & $d_a$ & $d_i$ (est.) & $\rho$ (est.) & ARI \\
\midrule
Wine$\dagger$ & Original            & 13  & 7.9  & 1.65 & 0.897 \\
Wine$\dagger$ & $+$10 noise dims    & 23  & 8.1  & 2.84 & 0.621 \\
Wine$\dagger$ & $+$20 noise dims    & 33  & 9.8  & 3.36 & 0.208 \\
Wine$\dagger$ & $+$50 noise dims    & 63  & 10.2 & 6.18 & 0.041 \\
\midrule
MNIST (PCA)  & $d_a = 784$ & 784 & 38.2 & 20.5 & 0.361 \\
MNIST (PCA)  & $d_a = 100$ & 100 & 35.1 & 2.85 & 0.374 \\
MNIST (PCA)  & $d_a = 20$  & 20  & 13.4 & 1.49 & 0.353 \\
\midrule
Olivetti     & Full ($d_a=4096$) & 4096 & 41.3 & 99.2 & 0.312 \\
Olivetti     & PCA ($d_a=100$)   & 100  & 38.7 & 2.58 & 0.318 \\
\bottomrule
\end{tabular}
\end{table}

The Wine results directly replicate the synthetic failure threshold.
The image datasets confirm an important caveat: $\rho$ is predictive of
algorithm failure specifically when cluster separability (under low noise) is
already reasonable --- the regime in which distance dilution is the binding
constraint.

% ── 6. Discussion ─────────────────────────────────────────────────────────────
\section{Discussion}

\paragraph{Practical implications.}
Our results suggest a simple diagnostic protocol: before running any clustering 
pipeline, estimate $\rho$ using TwoNN. If $\rho > 5$, dimensionality reduction 
is not optional --- the algorithm will likely fail regardless of parameter 
tuning. For kNN-based methods, high $\rho$ is less concerning and may even 
be beneficial.

\paragraph{The two-sided nature of $\rho$.}
The divergence between centroid-based and distance-weighted methods under 
increasing $\rho$ has a geometric interpretation. High $\rho$ implies that 
most ambient dimensions carry noise, which concentrates pairwise distances 
--- destroying cluster separation for k-means but creating more uniform local 
neighborhoods for kNN, reducing sensitivity to local density variations.

\paragraph{Limitations.}
Our scaling laws are empirical; theoretical proofs of the power law form and 
the logarithmic $\beta$ scaling remain open. The real-data validation used 
noise injection rather than naturally varying $\rho$, which may not fully 
replicate real-world conditions. Results may differ for deep representations 
where intrinsic dimension estimation is harder.

\paragraph{Future work.}
Natural extensions include: theoretical derivation of the scaling law from 
concentration of measure principles; extension to deep learning representations; 
characterization of $\rho^*$ for other algorithm families (SVM, Gaussian 
processes, spectral methods); and development of $\rho$-aware preprocessing 
pipelines.

% ── 7. Conclusion ─────────────────────────────────────────────────────────────
\section{Conclusion}

We introduced the stretching ratio $\rho = d_a / d_i$ as a unified predictor 
of machine learning algorithm behavior in high-dimensional settings. Through 
systematic experiments, we established a power law decay of clustering 
performance with $\rho$, a universal failure threshold at $\rho^* \approx 
3.5$--$5$, and a logarithmic scaling law for the decay exponent. We showed 
that $\rho$ has opposite effects on centroid-based versus distance-weighted 
methods, revealing a geometric tradeoff that transcends the simple curse of 
dimensionality narrative. Our framework provides practitioners with a single, 
interpretable diagnostic that predicts algorithm failure better than existing 
complexity measures individually.

\begin{thebibliography}{99}
\bibitem{bellman1957} Bellman, R.: Dynamic Programming. Princeton University Press (1957).
\bibitem{beyer1999} Beyer, K., Goldstein, J., Ramakrishnan, R., Shaft, U.: When is ``nearest neighbor'' meaningful? In: ICDT (1999).
\bibitem{bengio2013} Bengio, Y., Courville, A., Vincent, P.: Representation learning: A review and new perspectives. IEEE TPAMI 35(8), 1798--1828 (2013).
\bibitem{facco2017} Facco, E., d'Errico, M., Rodriguez, A., Laio, A.: Estimating the intrinsic dimension of datasets by a minimal neighborhood information. Scientific Reports 7(1), 12140 (2017).
\bibitem{levina2004} Levina, E., Bickel, P.: Maximum likelihood estimation of intrinsic dimension. In: NeurIPS (2004).
\bibitem{vershynin2018} Vershynin, R.: High-Dimensional Probability. Cambridge University Press (2018).
\bibitem{aggarwal2001} Aggarwal, C.C., Hinneburg, A., Keim, D.A.: On the surprising behavior of distance metrics in high dimensional space. In: ICDT (2001).
\end{thebibliography}

\section*{Figures}

% Note: If your image files aren't uploaded yet, these will show as blank boxes 
% with text but the project WILL compile perfectly without throwing errors.

\begin{figure}[h]
    \centering
    \includegraphics[width=0.9\textwidth]{fig1_main_result.png}
    \caption{(a) k-means ARI decreases monotonically with stretching ratio 
    $\rho$. Shaded regions show $\pm 1$ standard deviation across trials. 
    The red dotted line marks $\rho^*=5$. (b) Power law fits 
    ARI $\sim \alpha\rho^{-\beta}$ across all intrinsic dimensions.}
    \label{fig:main}
\end{figure}

\begin{figure}[h]
    \centering
    \includegraphics[width=0.9\textwidth]{fig2_algorithm_comparison.png}
    \caption{Algorithm-dependent effects of $\rho$. k-means degrades 
    (left) while kNN classification (center) and regression (right) improve. 
    This divergence reveals a geometric tradeoff rather than a simple curse.}
    \label{fig:algos}
\end{figure}

\begin{figure}[h]
    \centering
    \includegraphics[width=0.9\textwidth]{fig3_heavytail.png}
    \caption{Heavy-tailed noise shifts the failure threshold earlier. Under 
    $t$-distribution noise with df=2, failure occurs at $\rho^* < 2$ compared 
    to $\rho^* \approx 3$--$5$ under Gaussian noise.}
    \label{fig:heavytail}
\end{figure}

\begin{figure}[h]
    \centering
    \includegraphics[width=0.9\textwidth]{fig4_beta_and_predictors.png}
    \caption{(a) Decay exponent $\beta$ scales logarithmically with $d_i$. 
    (b) $\log(\rho)$ achieves highest $R^2$ among individual predictors, 
    outperforming participation ratio and explained variance ratio.}
    \label{fig:beta}
\end{figure}

\end{document}
